\section{Lecture 7, Part 0: Approximate Nearest Neighbour Search}

\begin{frame}
\vfill
\centering
{\color{RLBlue}\large Lecture 7}

\vspace{0.3em}
{\color{RLBlue}\Large\bfseries Part 0}

\vspace{0.6em}
{\Large Approximate Nearest Neighbour Search}
\vfill
\end{frame}

\begin{frame}{Learning Goals}
\begin{itemize}
  \item Say when an approximate index is needed at all, and compute the corpus size at which exact search stops being fast enough.
  \pause
  \item Explain how IVF, product quantization and HNSW each cut the work, and name exactly where each one loses recall.
  \pause
  \item Derive why HNSW's layer assignment is geometric, and read a recall-versus-work curve as an engineering budget.
\end{itemize}
\end{frame}

% ─── The scale wall ─────────────────────────────────────────────────────────
\begin{frame}{Why Is Exact Search Too Slow?}
\begin{columns}[T]
\column{0.50\textwidth}
\small
\begin{itemize}
  \item Lecture 6 turned every article into a vector. Finding the closest one is a dot product against \emph{every} stored vector -- $N$ of them, each $D$-dimensional.
  \pause
  \item That is $O(ND)$ per query, and it is entirely predictable. At $D = 1024$: about $2$\,ms at $7{,}368$ vectors, $30$\,ms at $10^5$, close to $300$\,ms at $10^6$.
  \pause
  \item \textbf{Nothing is wrong with the code.} The work is linear in $N$, so a corpus $100\times$ larger costs $100\times$ more -- there is no constant factor left to win.
\end{itemize}
\column{0.48\textwidth}
\centering
\includegraphics[width=\linewidth]{figures/ann/ann_scale_wall.pdf}

\vspace{0.2em}
\scriptsize Measured on one laptop core, \texttt{float32}, $D = 1024$.
\end{columns}
\end{frame}

\begin{frame}{Worked Example: Our Corpus Does Not Need an Index}
\begin{itemize}
  \item $7{,}368$ articles at $D = 1024$ in \texttt{float32} is $7368 \times 1024 \times 4 = 30.2$\,MB. It fits in cache-friendly RAM on any laptop.
  \pause
  \item A brute-force scan answers in roughly $2$\,ms with \textbf{recall $= 1.000$ by construction} -- there is no approximation, so there is nothing to lose.
  \pause
  \item An HNSW index over the same $7{,}368$ vectors might reach $0.3$\,ms. You would have bought $1.7$\,ms and paid for it with a build step, a dependency, and a recall below $1$.
  \pause
  \item \textbf{Rule of thumb:} below roughly $10^5$ vectors, scan exhaustively. This lecture is about $10^6$ to $10^9$, where that stops being an option.
\end{itemize}
\pause
\begin{rlpitfall}
Reaching for HNSW on a $7{,}000$-vector corpus adds a recall loss for no latency gain. Knowing when \emph{not} to index is part of knowing how to index.
\end{rlpitfall}
\end{frame}

\begin{frame}{Nearest Neighbours, Exactly and Approximately}
\begin{rldefinition}
\textbf{Exact $k$-NN} -- return the $k$ stored vectors with the highest similarity to the query. \emph{Here:} the $10$ articles whose embeddings are closest to ``maximum working hours per week''.

\smallskip
\textbf{Approximate $k$-NN (ANN)} -- return $k$ vectors that are \emph{probably} those, having examined far fewer than $N$ candidates.

\smallskip
\textbf{Recall@$k$} -- of the $k$ true nearest neighbours, the fraction the index actually returned. \emph{Here:} recall@$10 = 0.9$ means nine of the ten correct articles came back and one was missed.
\end{rldefinition}
\pause
\begin{itemize}
  \item Every index in this lecture trades among the same three resources: \textbf{recall}, \textbf{latency}, and \textbf{memory}. You may pick two.
\end{itemize}
\end{frame}

% ─── IVF ────────────────────────────────────────────────────────────────────
\begin{frame}{IVF: Only Look in the Nearby Cells}
\small
Cluster the corpus once with $k$-means into \texttt{nlist} cells. At query time, find the nearest centroids and scan \emph{only} the vectors filed under them.
\vspace{0.2em}

\centering
\begin{tikzpicture}[
  font=\scriptsize,
  cellbox/.style={draw=RLBlue!55, line width=0.5pt},
  probed/.style={draw=RLBlue, line width=0.9pt, fill=RLOrange!35},
  cen/.style={circle, fill=RLBlue, inner sep=1.4pt},
  pt/.style={circle, fill=RLGreen!85, inner sep=0.7pt},
  lbl/.style={font=\scriptsize, inner sep=1.5pt}
]
  % probed cells shaded first, so grid lines draw on top
  \fill[RLOrange!40] (2.1,0) rectangle (3.95,1.7);
  \fill[RLOrange!16] (2.1,1.7) rectangle (3.95,3.3);
  % outer box and cell boundaries
  \draw[cellbox] (0,0) rectangle (6,3.3);
  \draw[cellbox] (2.1,0) -- (2.1,3.3);
  \draw[cellbox] (3.95,0) -- (3.95,3.3);
  \draw[cellbox] (0,1.7) -- (6,1.7);
  % centroids
  \node[cen] at (1.05,2.5) {}; \node[cen] at (3.02,2.5) {}; \node[cen] at (4.95,2.5) {};
  \node[cen] at (1.05,0.85) {}; \node[cen] at (3.02,0.85) {}; \node[cen] at (4.95,0.85) {};
  % scattered vectors
  \foreach \x/\y in {0.5/2.9, 1.4/2.8, 0.8/2.1, 1.6/2.3, 0.6/1.2, 1.5/0.5, 1.7/1.3,
                     2.5/3.0, 3.6/2.9, 2.7/2.0, 3.6/2.2, 2.4/1.2, 3.5/0.5, 3.0/1.4,
                     4.4/2.9, 5.4/2.7, 4.6/2.0, 5.5/2.2, 4.5/1.1, 5.4/0.6, 5.0/1.5}
    \node[pt] at (\x,\y) {};
  % query
  \node[circle, fill=red!75!black, inner sep=1.7pt] (q) at (3.25,1.05) {};
  \node[lbl, anchor=west] at (3.42,1.05) {query};
  \node[lbl, anchor=south west, font=\tiny] at (2.15,0.03) {probed};
  \node[lbl, anchor=south west, font=\tiny] at (2.15,1.73) {2nd probe};
\end{tikzpicture}

\vspace{0.1em}
\footnotesize
\texttt{nprobe} sets how many cells are opened. Recall is lost exactly when a true neighbour sits in a cell you did not open.
\end{frame}

\begin{frame}{Worked Example: What IVF Buys and Costs}
\begin{itemize}
  \item A common heuristic is $\texttt{nlist} \approx \sqrt{N}$, so each cell holds about $\sqrt{N}$ vectors. At $N = 10^6$: $1{,}000$ cells of $1{,}000$ vectors each.
  \pause
  \item Query cost is the coarse step plus the scan: $\texttt{nlist} + \texttt{nprobe}\times\frac{N}{\texttt{nlist}}$ distance computations. With $\texttt{nprobe} = 8$ that is $1{,}000 + 8{,}000 = 9{,}000$ instead of $1{,}000{,}000$ -- about a $110\times$ saving.
  \pause
  \item \textbf{Memory cost: none.} IVF stores the same vectors, merely grouped, plus \texttt{nlist} centroids. It buys latency with recall alone.
  \pause
  \item \textbf{Where recall goes:} a true neighbour near a cell boundary is filed on the far side. Raising \texttt{nprobe} recovers it and gives back the saving.
\end{itemize}
\end{frame}

% ─── PQ ─────────────────────────────────────────────────────────────────────
\begin{frame}{Product Quantization: A Vector in $64$ Bytes}
\small
Split each vector into $M$ subvectors. Run $k$-means with $256$ centroids \emph{inside each subspace}, and store only which centroid each piece landed on -- one byte per piece.
\vspace{0.3em}

\centering
\begin{tikzpicture}[
  font=\scriptsize,
  seg/.style={draw=RLBlue, fill=RLGray, minimum width=13mm, minimum height=5mm, inner sep=1pt},
  code/.style={draw=RLGreen, fill=RLOrange!25, minimum width=13mm, minimum height=5mm, inner sep=1pt},
  book/.style={draw=RLGreen!70, fill=RLGray, minimum width=13mm, minimum height=4mm, inner sep=1pt, font=\tiny},
  lbl/.style={font=\scriptsize, inner sep=1.5pt}
]
  \node[lbl, anchor=east] at (-0.15,0) {one vector, $D=1024$};
  \node[seg] (s1) at (0.8,0) {dims 1--16};
  \node[seg] (s2) at (2.25,0) {17--32};
  \node[seg] (s3) at (3.7,0) {33--48};
  \node[lbl] (dots) at (4.95,0) {$\cdots$};
  \node[seg] (s4) at (6.2,0) {1009--1024};

  \node[book] (b1) at (0.8,-1.15) {codebook 1};
  \node[book] (b2) at (2.25,-1.15) {codebook 2};
  \node[book] (b3) at (3.7,-1.15) {codebook 3};
  \node[book] (b4) at (6.2,-1.15) {codebook 64};
  \foreach \i in {1,2,3,4} \draw[->, RLGreen] (s\i) -- (b\i);

  \node[code] (c1) at (0.8,-2.3) {\texttt{0x2A}};
  \node[code] (c2) at (2.25,-2.3) {\texttt{0xB1}};
  \node[code] (c3) at (3.7,-2.3) {\texttt{0x07}};
  \node[lbl] at (4.95,-2.3) {$\cdots$};
  \node[code] (c4) at (6.2,-2.3) {\texttt{0x9C}};
  \foreach \i in {1,2,3,4} \draw[->, RLGreen] (b\i) -- (c\i);

  \node[lbl, anchor=west] at (7.05,-2.3) {$64$ bytes total};
  \node[lbl, anchor=west] at (7.05,0) {$4096$ bytes};
\end{tikzpicture}

\vspace{0.2em}
\footnotesize
Each codebook holds $256$ centroids, so one byte names a piece. $M = 64$ pieces $\Rightarrow$ $64$ bytes per vector.
\end{frame}

\begin{frame}{Worked Example: What PQ Costs and Saves}
\begin{itemize}
  \item \textbf{Memory.} A \texttt{float32} vector at $D = 1024$ is $D \times 4 = 4096$ bytes. With $M = 64$ it becomes $64$ bytes -- a $64\times$ reduction. At $10^9$ vectors that is $4096$\,GB down to $64$\,GB.
  \pause
  \item \textbf{The codebooks are free.} $64$ subspaces $\times\ 256$ centroids $\times\ 16$ dims $\times\ 4$ bytes $= 1$\,MiB, once, for the whole corpus.
  \pause
  \item \textbf{Speed, via a lookup table.} Compute the query's similarity to all $256$ centroids in each subspace up front: a $64 \times 256$ table of $64$\,KB, small enough to stay in cache. Every stored vector then costs $64$ lookups and adds -- no multiplications at all.
  \pause
  \item \textbf{Where recall goes:} the stored vector \emph{is} its centroid now. Two articles quantized to the same code are indistinguishable, and the error is permanent.
\end{itemize}
\end{frame}

\begin{frame}{The Recall/Work Trade-Off, Measured}
\begin{columns}[T]
\column{0.48\textwidth}
\small
\begin{itemize}
  \item IVF at $\texttt{nprobe} = 1$ scans $0.8\%$ of the corpus and recovers only $30\%$ of the true top ten. At $\texttt{nprobe} = 16$ it scans $11\%$ and recovers $94\%$.
  \pause
  \item PQ alone is worse than either -- at $32\times$ compression it recovers about a quarter. \textbf{So nobody uses PQ alone:} it produces a \emph{shortlist}, which is then re-scored with the exact vectors.
  \pause
  \item Re-ranking $2.5\%$ of the corpus exactly recovers $94\%$, at a thirty-second of the memory.
\end{itemize}
\column{0.50\textwidth}
\centering
\includegraphics[width=\linewidth]{figures/ann/ann_recall_tradeoff.pdf}

\vspace{0.2em}
\scriptsize Measured on $20{,}000$ synthetic clustered vectors, $300$ queries, against exhaustive ground truth. The \emph{shape} transfers; absolute recall on real embeddings will differ.
\end{columns}
\end{frame}

\begin{frame}{Reading That Curve Honestly}
\begin{itemize}
  \item The two curves are not a fair race. PQ's cheap pass still touches \emph{every} vector -- it is only cheap per vector. IVF skips most vectors entirely but stores them at full width.
  \pause
  \item So they trade different resources: \textbf{PQ spends recall to save memory}; \textbf{IVF spends recall to save computation}. The $x$-axis counts only the exact distance computations, which is why PQ looks cheaper than it is.
  \pause
  \item In practice the two compose -- \texttt{IVFPQ} files PQ codes into IVF cells, so you skip most cells \emph{and} store the survivors compressed. FAISS (Johnson et al., 2019) is where most people first meet the combination.
\end{itemize}
\end{frame}

% ─── HNSW ───────────────────────────────────────────────────────────────────
\begin{frame}{HNSW: A Skip List Over a Proximity Graph}
\small
Build a graph whose edges join near neighbours, then stack sparser copies above it. Search enters at the top, walks greedily downhill, and drops a layer when it can get no closer.
\vspace{0.2em}

\centering
\begin{tikzpicture}[
  font=\scriptsize,
  nd/.style={circle, draw=RLBlue, fill=RLGray, inner sep=1.1pt},
  hot/.style={circle, draw=RLBlue, fill=RLOrange, inner sep=1.4pt},
  lnk/.style={draw=RLGreen!70, line width=0.4pt},
  lbl/.style={font=\scriptsize, inner sep=1.5pt}
]
  % layer 2
  \node[lbl, anchor=east] at (-0.2,3.4) {layer $2$};
  \node[hot] (e) at (1.0,3.4) {};
  \node[nd]  (a2) at (4.2,3.4) {};
  \draw[lnk] (e) -- (a2);
  \node[lbl, anchor=west] at (4.5,3.4) {$\sim N/M^2$ nodes};

  % layer 1
  \node[lbl, anchor=east] at (-0.2,1.9) {layer $1$};
  \node[hot] (b1) at (1.0,1.9) {};
  \node[nd] (b2) at (2.4,1.9) {};
  \node[hot] (b3) at (3.5,1.9) {};
  \node[nd] (b4) at (5.2,1.9) {};
  \draw[lnk] (b1) -- (b2) -- (b3) -- (b4);
  \draw[RLBlue, line width=0.9pt, ->] (b1) to[bend left=22] (b3);
  \node[lbl, anchor=south, RLBlue] at (2.25,2.35) {walk};
  \node[lbl, anchor=west] at (5.5,1.9) {$\sim N/M$ nodes};

  % layer 0
  \node[lbl, anchor=east] at (-0.2,0.4) {layer $0$};
  \foreach \i/\x in {1/0.6, 2/1.4, 3/2.2, 4/2.9, 6/4.3, 7/5.0, 8/5.8}
    \node[nd] (c\i) at (\x,0.4) {};
  \node[hot] (t) at (3.6,0.4) {};
  \draw[lnk] (c1) -- (c2) -- (c3) -- (c4) -- (t) -- (c6) -- (c7) -- (c8);
  \draw[lnk] (c2) to[bend right=28] (c4);
  \draw[lnk] (c4) to[bend right=28] (c6);
  \node[lbl, anchor=west] at (6.1,0.4) {all $N$ nodes};

  % descent
  \draw[->, RLBlue, dashed] (e) -- (b1);
  \draw[->, RLBlue, dashed] (b3) -- (t);
  \node[lbl, anchor=east, RLBlue] at (0.9,2.65) {descend};
  \node[lbl, anchor=west, RLBlue] at (3.65,1.15) {descend};
\end{tikzpicture}

\vspace{0.1em}
\footnotesize
Orange marks the path: enter at the top, move to the closest neighbour until none is closer, then drop.
\end{frame}

\begin{frame}{Why the Layer Assignment Is Geometric}
\begin{itemize}
  \item Each inserted node is given a maximum level $\ell = \lfloor -\ln(U)\, m_L \rfloor$ with $U \sim \mathrm{Uniform}(0,1)$. Then
  \vspace{-0.4em}
  \[
    P(\ell \ge j) \;=\; P\!\left(U \le e^{-j/m_L}\right) \;=\; e^{-j/m_L}.
  \]
  \vspace{-1.2em}
  \pause
  \item Choosing $m_L = 1/\ln M$ makes this exactly $M^{-j}$: \textbf{each layer up keeps one node in $M$}. Layer sizes are $N,\; N/M,\; N/M^2, \dots$
  \pause
  \item So the top layer is reached when $N M^{-j} \approx 1$, i.e.\ $j \approx \log_M N$. At $M = 16$, $N = 10^6$: $\log_{16}(10^6) = 4.98$, about five layers, with $62{,}500$ nodes on layer~1.
  \pause
  \item That is precisely a skip list: each layer is an express lane with a fraction $1/M$ of the stops, so the descent costs $O(\log N)$ hops instead of $O(N)$ comparisons (Malkov \& Yashunin, 2018).
\end{itemize}
\end{frame}

\begin{frame}{What HNSW Costs}
\begin{itemize}
  \item \textbf{Memory.} Layer $0$ allows up to $2M$ links per node, the layers above $M$; the geometric tail adds almost nothing. At $M = 16$ that is about $33$ links per node, $4$ bytes each -- roughly $132$ bytes.
  \pause
  \item On $10^6$ vectors of $D = 1024$: the graph costs $0.13$\,GB on top of $4.10$\,GB of vectors -- about $3\%$ overhead. HNSW is cheap in memory but does \emph{not} compress; that is what PQ is for.
  \pause
  \item \textbf{Search is tuned by \texttt{efSearch}}, the size of the candidate list kept at layer $0$. Larger means more of the graph explored, higher recall, more time -- the same dial as IVF's \texttt{nprobe}.
  \pause
  \item \textbf{Where recall goes:} greedy descent can settle in a local optimum and never reach the true neighbour's neighbourhood. A wider candidate list makes that rarer, never impossible.
\end{itemize}
\end{frame}

\begin{frame}{ScaNN: Quantization Aimed at Ranking}
\begin{itemize}
  \item Plain PQ minimises \emph{reconstruction} error -- how far the centroid sits from the original vector. But we never wanted the vector back; we wanted its \emph{ranking} against a query.
  \pause
  \item Guo et al.\ (2020) observe that those are not the same objective. For maximum inner-product search, error \emph{along} the datapoint's direction distorts the score much more than error perpendicular to it.
  \pause
  \item Their \textbf{anisotropic} quantization weights the two directions differently when fitting the codebook, buying accuracy on exactly the comparisons that decide the top-$k$.
  \pause
  \item \textbf{The general lesson, well beyond ScaNN:} optimise the quantity you are actually scored on. It is the same move as Lecture 8's insistence on scoring trajectories rather than final answers.
\end{itemize}
\end{frame}

\begin{frame}{Where Each Method Loses Recall}
\centering
\small
\begin{tabular}{@{}l|l|l@{}}
\textbf{Method} & \textbf{Saves} & \textbf{Loses a neighbour when\ldots} \\
\hline
IVF & computation & it sits in an unprobed cell \\[2pt]
PQ & memory & its code collides with another vector's \\[2pt]
HNSW & computation & greedy descent stalls before reaching it \\[2pt]
IVFPQ & both & either of the first two happens \\
\end{tabular}

\vspace{0.8em}
\begin{itemize}
  \item \textbf{None of them can tell you what they missed.} Recall is measurable only against exhaustive ground truth, which is why you compute it once on a sample rather than trusting a default.
  \pause
  \item That single habit -- build the exact answer for a few hundred queries, then tune the index against it -- is the whole of index selection.
\end{itemize}
\end{frame}
